Tijdens haar abductie
had ze een reflectie.

Dat was vaak een zegen,
maar werkte soms tegen

als het leidde tot ‘ik
kijk naar de ander'-blik.

De illusie van iets
beziet een ander niets.

Geen zelfobservatie,
maar een zelfcorrectie

op grond van een oordeel.
Denkwerk werd een nadeel

en zelfs een lijntrekker,
die volgens Heidegger

slechts de wereld vervlakt,
als het alleen abstract

het authentiek bestaan
vervreemd durft aan te gaan.

Reflectie werd aldus
een soort mentale lus,

waarin je observeert
en diep analyseert,

maar er werd niet doorvoeld
wat er diep in je woelt.

Al eerder gegrepen
had ze ooit begrepen

dat ze kon opzwepen
door slimme ingrepen.

Die eerdere weiman
was nog een jongeman.

Ze ontstak al een lont
als ze dicht voor hem stond.

Hij raakte ook bij haar
een gevoelige snaar

als zijnde een droomprins.
Dat was vreemd enigszins.

Vele premiejagers
waren vaak baarddragers.

Ze oogden vrij lelijk
en werkten bedrieglijk.

De vaak ex-bandieten
konden erg genieten

van het leed van hun prooi.
Ze vonden kwellen mooi.

Emmy keek op hen neer.
Dat gold niet voor haar heer.

Hij was gemakkelijk
en zeer aantrekkelijk.

Hij bleef vrij onschuldig
en wachtte geduldig

toen hij haar destijds ving
op haar toenadering.

Zijn aura was roze
in mistige, losse

en zachte patronen
die zijn lichaam kronen.

Hij was ook ondeugend.
Dat vond ze verheugend.

Evenals de sculptuur
van zijn sterke postuur.

Slanke lange benen
die haar entertainen.

Zware lange armen
die konden verwarmen.

Een zonnetje dat scheen
door zijn goudlokken heen.

Qua bekoring als sfeer
behoorde de jonkheer

wis bij de spekkopers.
Johannes de Doper’s

grote vreugdevuren
smeulden al wat uren,

toen Emmy werd betrapt
en behendig gesnapt.

Vlamde de vernieuwing?
Wachtte haar zuivering?

Nadat ze was geboeid
spraken ze onvermoeid

lopende de voettocht,
waarbij zij rusten mocht

zo vaak als ze wilde.
Soms als hij wild grilde,

boven een heet kampvuur,
bekeek hij haar figuur.

Op dat punt geen domkop
ving zij die blikken op.

Haar vogelvrij tenue
nodigde continu

uit tot stiekem gluren.
Ze kon het verduren.

Haar ooit golvende rok
hing nu kaduuk halfstok.

Door wind afgeraffeld
en regen verstaffeld

hing het halverwege
Emmy’s hooggelegen

smalle bovenbenen.
Door de zon beschenen.

Een jurk of dameshoed
stond op gespannen voet

met vluchten en jagen.
Geen schelm zal die dragen.

Het onstuimig leven
om te overleven

door de vogelvrijen
scheurde bij de dijen

alle minirokjes.
Ze hielden met stokjes

de stukken bij elkaar.
Daardoor toonde bij haar

in de beide spleten
de bijna versleten

onderrok die ze droeg.
De weiman zag genoeg.

Borstjes werden bedekt
door een stof uitgerekt

tot over, als houder,
Emmy’s linkerschouder.

Deze bustewonder
liep gespannen onder
haar rechter oksel door
en ze knoopte ervoor

de uiteinden samen
die scheef bijeenkwamen.

Toch zakte een tietje
uit haar bikinietje

zichtbaar van tijd tot tijd.
Van dat had ze geen spijt.

Als ze weer werd bespied
en voldeed duwen niet,

dan knoopte ze het doek
in een bepaalde hoek,

zodat hij goed kon zien
wat de gehaaide trien

onder haar hoede had.
De weiman vond het wat.

Haar positie wetend,
maar haar staat vergetend,

wou ze hem ophitsen
om hem snel te gidsen

naar een geilheid
verward met verliefdheid.

Dus droogde op een stok
haar dunne onderrok.

Gewassen in een beek,
terwijl hij haar bekeek.

Zo hing slechts een lapje
voor haar rode kapje.

Die toonde zich spontaan
als ze gebukt ging staan.

Hij wierp haar zijn blik toe.
Ze speelde kiekeboe

met haar flamoes of tiet.
Dat beeld weerstond hij niet.

Het prachtige gezicht
van het rossige wicht

met haar hemels lichaam
strikte hem te langzaam.

Dus koos slimme Emmy
als slinkse strategie

onthullende poses
als een extra dosis.

Eenmaal buiten zinnen
wist ze hem te winnen.

De jonkheer zei al gauw:
“Ik ben verliefd op jou!

Ik wil je niet meesjouwen,
maar graag met je trouwen.”.

Zonder een aarzeling
stemde de vluchteling

met zijn verloving in.
Ze had immers haar zin.

Doch werd hij ook strafbaar.
Dat verontrustte haar.

Hij negeerde dat feit
door zijn pure geilheid.

Thomas van Aquino
zei het in zijn pleidooi

zo treffend en zo raak:
“Wellust benevelt vaak

de rede van de mens”.
Zijn drift was te intens.

Hij maakte Emmy los
en stalde haar op mos.

Onder de paarse lucht
slaakten beiden een zucht

naar de zonsondergang.
Emmy voelde geen dwang.

Ze draaide naar hem toe.
“God zal wel weten hoe,

maar ik wil deze prins.
Geef me enkele hints.

Is hij wel de ware?
Moet ik me nog sparen?”

Alsof hij als spion
gedachten lezen kon,

streelde hij haar gezicht,
die al lang was gezwicht.

De vorst gaf haar een kus.
Van haar oudere zus

wist ze wat ze moest doen
bij haar eerste tongzoen.

Toen zijn persende tong
door een spleet binnendrong

doolde die even rond
door Emmy’s open mond

als slurpende zuigsnuit.
Het kampvuur doofde uit.

Het werd pikkedonker.
Ze dacht; “pik de jonker".

“Slik!”, dacht de jongeheer.
Dik werd zijn jongeheer.

Haar rode hoofd tolde.
Haar natte tong rolde

vief over die van hem.
Hij hield haar in een klem.

Ze streelde zijn lokken.
“Gaat hij mij nu fokken?”

Emmy rilde van angst.
In het duister het bangst.

Nergens was een luipaard
die haar had aangestaard

en zichzelf had verlinkt.
Toch gaf Emmy’ instinct,

dezelfde reactie,
als was ze een bambi.

Was de mens het roofdier
met haar in het vizier

dat het duister aangreep
voor een blinddaterape?
